\documentclass{zhslides}

\usepackage{logic}


\title{Day 1: 哲学的语言转向与逻辑学基础}
%\subtitle{}
\date{\today}
\date{2026年7月20日}
\author{于佳铭}
\institute{唯理书院2026\\精确的逻辑与模糊的意义}

\begin{document}

\maketitle

%\begin{frame}{目录}
%  \tableofcontents
%\end{frame}

% 可以用 \alert{} 来高亮



\section{语言为什么是个问题？}

\begin{frame}{语言为什么是个问题？}

  \textbf{语言富有歧义}
  \pause
  
  \vspace{1em}
  \begin{quote}
    我买了橙子或苹果。
  \end{quote}
  \vspace{1em}\pause
  \begin{quote}
    每个孩子都喜欢一个玩具。
  \end{quote}
%  \begin{quote}
%    每个学生都没有通过某门考试。
%  \end{quote}
  \vspace{1em}\pause
  \begin{quote}
    每个人都在享受着好时光。
  \end{quote}

\end{frame}



\begin{frame}{语言为什么是个问题？}
  
  \textbf{语言是不精确、混乱的}
  \pause
  
  \vspace{1em}
  \begin{quote}
    天空和草地是同一种颜色。
  \end{quote}
  \pause
  \bit
    \item 中文里区分“蓝”“绿”，而古日语里只用“青”一种颜色\footnote<.->[frame]{\url{https://en.wikipedia.org/wiki/Blue\%E2\%80\%93green_distinction_in_language}}
    \item 语言在描述世界的同时也在构建我们对世界的理解
  \eit

\end{frame}


\begin{frame}[t]{哲学的语言转向}

  \vspace{2em}
  \begin{brainstorm*}{}
    为什么语言的不精确对于哲学可能是个问题？
  \end{brainstorm*}

  \pause
  \only<2>{
  \bit
    \item 传统哲学：心灵是什么？自由意志存在吗？道德真理是客观的吗？
  \eit
  }
  
  \bit
    \item 早期维特根斯坦（《逻辑哲学论》）：语言图像论
      \bit
        \item 任何有意义的命题对应世界的某个状态（\term{逻辑图像}）
        \item 乐谱——交响乐
        \item 哲学问题可以通过形式化语言迎刃而解
        \item \textit{“我语言的界限就是我世界的界限。”}
      \eit
      \vspace{0.5em}
    \item 罗素（\textit{Our Knowledge of the External World}）
      \bit
        \item 自然语言阻碍我们认识到现实世界的\term{逻辑形式}{logical form}
        \item 哲学的困惑源自语言的逻辑与现实经验的不对等
      \eit
  \eit

\end{frame}


\begin{frame}[t]{哲学的语言转向}

  \vspace{2em}
  \begin{brainstorm*}{}
    为什么语言的不精确对于哲学可能是个问题？
  \end{brainstorm*}

\bit
  \item 晚期维特根斯坦（《哲学研究》）：语言游戏论
    \bit
      \item 意义即图像 \raisebox{0.2ex}{$\longrightarrow$} 意义即使用
    \eit
\eit

\end{frame}






\section{逻辑学基础}

\begin{frame}{逻辑学基础}
  \bit
    \item 经典逻辑的研究对象和基本单位：\term{命题}{proposition}
    \item 命题是能够判断真假（记作$1, 0$或$\true, \false\hspace{0.1em}$）的陈述句
    \item \term{形式化语言}{formal language}是用精确的公式与规则定义的语言
    \item 我们用字母表示命题：$P, Q, R, P_1, P_2, \dots$
  \eit
\end{frame}


\begin{frame}[t]{逻辑学基础}
  
  \vspace{2em}
  \begin{brainstorm*}{}
    为什么我们想要一个形式化语言？
  \end{brainstorm*}
  \vspace{1em}
  
  \begin{quote}
    如果下雨，我就带伞。我带了伞。所以，下雨了。
  \end{quote}
  \vspace{1em}\pause
  \begin{quote}
    猫都是动物。狗都是动物。所以，猫都是狗。
  \end{quote}
  \vspace{1em}\pause
  \begin{quote}
    没有狗能飞。没有猫是狗。所以，猫都能飞。
  \end{quote}

\end{frame}


\begin{frame}{命题逻辑符号}
  
  \begin{table}[h]
  \centering
  {\renewcommand{\arraystretch}{1.2}
  \begin{tabular}{llc}
    \textbf{自然语言对应} & \textbf{名称} & \textbf{逻辑符号} \\\hline
    不 & 逻辑非 & $\neg$ \\\pause
    并且 & 逻辑与 & $\land$ \\\pause
    或者 & 逻辑或 & $\lor$ \\\pause
    如果……那么…… & 材料蕴含 & $\implies$
  \end{tabular}
  }
  \end{table}
  
\end{frame}


\begin{frame}{真值表}

  \begingroup
  \setlength{\tabcolsep}{4pt}
  \begin{table}[h]
  \centering
  \begin{tabular}{c|c}
    $\neg$ & \\\hline
    1 & 0 \\
    0 & 1
  \end{tabular}
  \qquad
  \begin{tabular}{c|cc}
    $\land$ & 1 & 0 \\\hline
    1 & 1 & 0 \\
    0 & 0 & 0
  \end{tabular}
  \qquad
  \begin{tabular}{c|cc}
    $\lor$ & 1 & 0 \\\hline
    1 & 1 & 1 \\
    0 & 1 & 0
  \end{tabular}
  \qquad
  \begin{tabular}{c|cc}
    $\implies$ & 1 & 0 \\\hline
    1 & 1 & 0 \\
    \only<1>{0 & 1 & 1}
%    \only<2>{\rowcolor{yellow} 0 & 1 & 1}
  \end{tabular}
  \end{table}
  \endgroup

\end{frame}


\begin{frame}{逻辑形式化}

  \begin{quote}
    如果下雨，我就带伞。我带了伞。所以，下雨了。
  \end{quote}
  \vspace{0.5em}\pause
  \begin{quote}
    $\quad R \implies S,\ S$ \\
    $\quad \therefore R$
  \end{quote}
  
  \vspace{1em}\pause
  \begin{quote}
    猫都是动物。狗都是动物。所以，猫都是狗。
  \end{quote}
  \vspace{1em}\pause
  \begin{quote}
    没有狗能飞。没有猫是狗。所以，猫都能飞。
  \end{quote}

\end{frame}


\begin{frame}{一阶逻辑符号}
  
  \begin{table}[h]
  \centering
  {\renewcommand{\arraystretch}{1.2}
  \begin{tabular}{llc}
    \textbf{自然语言对应} & \textbf{名称} & \textbf{逻辑符号} \\\hline
    存在 & 存在量词 & $\exists$ \\\pause
    任意/所有 & 全称量词 & $\forall$ \\
  \end{tabular}
  }
  \end{table}
  
\end{frame}


\begin{frame}{一阶逻辑元素}
  
  \bit
    \item \term{谓词}{predicate}：$P, Q, R, \dots$
    \item \term{常量}{constant}：$a, b, c, \dots$
    \item \term{变量}{variable}：$x, y, z, \dots$
  \eit
  \pause
  我们会使用希腊字母$\phi, \psi, \dots$来指代完整的一阶逻辑句子。
  
\end{frame}


\begin{frame}{一阶逻辑形式化}

  \begin{quote}
    一个男人抽烟。
    \pause\hfill
    $\exists x \left( \pred{Man}(x) \land \pred{Smokes}(x) \right)$
  \end{quote}
  
  \vspace{1em}\pause
  \begin{quote}
    只要是孩子就有喜欢的玩具。
  \end{quote}
  
  \vspace{1em}\pause
  \begin{quote}
    每个孩子都喜欢一个玩具。
  \end{quote}

  \vspace{1em}\pause
  \begin{quote}
    每个人都在享受着好时光。
  \end{quote}

  \vspace{1em}\pause
  \begin{quote}
    Every farmer who owns a donkey beats it.
  \end{quote}

\end{frame}












\end{document}